% Preámbulo modular para presentaciones Beamer (Modelación Estadística)
% Diseño y paleta institucional del Tecnológico de Monterrey

\documentclass[aspectratio=169,xcolor={dvipsnames,table}]{beamer}

\usepackage[utf8]{inputenc}
\usepackage[spanish,mexico]{babel}
\usepackage{amsmath,amssymb,amsfonts,mathtools}
\usepackage{graphicx}
\usepackage{listings}
\usepackage{booktabs}
\usepackage{tikz}
\usetikzlibrary{matrix,arrows,decorations.pathmorphing,shapes.geometric,positioning}

% Paleta Institucional Tecnológico de Monterrey
\definecolor{TecAzul}{HTML}{0039A6}       % Azul TEC: color primario institucional
\definecolor{TecAzulOscuro}{HTML}{1A2E51} % Azul marino oscuro
\definecolor{TecRojo}{HTML}{EC2661}       % Rojo de acento (paleta miTec)
\definecolor{TecGris}{HTML}{646464}       % Gris neutro para comentarios y subtítulos
\definecolor{TecFondoBloque}{HTML}{F4F6F9}% Fondo suave para bloques

% Configuración de Tema Beamer
\usetheme{Boadilla}
\usecolortheme[named=TecAzulOscuro]{structure}
\setbeamercolor{alerted text}{fg=TecRojo}
\setbeamercolor{palette primary}{bg=TecAzulOscuro,fg=white}
\setbeamercolor{palette secondary}{bg=TecAzul,fg=white}
\setbeamercolor{palette tertiary}{bg=TecAzulOscuro!80!black,fg=white}
\setbeamercolor{title}{fg=TecAzulOscuro,bg=TecFondoBloque}
\setbeamercolor{frametitle}{fg=TecAzulOscuro,bg=TecFondoBloque}

% Bloques personalizados elegantes
\setbeamercolor{block title}{bg=TecAzulOscuro,fg=white}
\setbeamercolor{block body}{bg=TecFondoBloque,fg=black}
\setbeamercolor{block title alerted}{bg=TecRojo,fg=white}
\setbeamercolor{block body alerted}{bg=TecRojo!8,fg=black}
\setbeamercolor{block title example}{bg=TecAzul,fg=white}
\setbeamercolor{block body example}{bg=TecAzul!8,fg=black}

% Redondear bordes y añadir sombra sutil
\setbeamertemplate{blocks}[rounded][shadow=true]
\setbeamertemplate{navigation symbols}{}

% Pie de página limpio con número de diapositiva
\setbeamertemplate{footline}{
  \leavevmode%
  \hbox{%
  \begin{beamercolorbox}[wd=.7\paperwidth,ht=2.25ex,dp=1ex,left]{author in head/foot}%
    \hspace*{2ex}\insertshortauthor~~(\insertshortinstitute) \hfill \insertshorttitle
  \end{beamercolorbox}%
  \begin{beamercolorbox}[wd=.3\paperwidth,ht=2.25ex,dp=1ex,right]{date in head/foot}%
    \insertshortdate{}\hspace*{2em}
    \insertframenumber{} / \inserttotalframenumber\hspace*{2ex}
  \end{beamercolorbox}}%
  \vskip0pt%
}

% Configuración de Listings de Python para visualización pedagógica de código
\lstset{
    language=Python,
    basicstyle=\ttfamily\footnotesize,
    keywordstyle=\color{TecAzulOscuro}\bfseries,
    stringstyle=\color{TecRojo},
    commentstyle=\color{TecGris}\itshape,
    showstringspaces=false,
    breaklines=true,
    frame=lines,
    rulecolor=\color{TecAzulOscuro!30},
    backgroundcolor=\color{TecFondoBloque},
    aboveskip=0.5em,
    belowskip=0.5em,
    literate={á}{{\'a}}1 {é}{{\'e}}1 {í}{{\'i}}1 {ó}{{\'o}}1 {ú}{{\'u}}1
             {Á}{{\'A}}1 {É}{{\'E}}1 {Í}{{\'I}}1 {Ó}{{\'O}}1 {Ú}{{\'U}}1
             {ñ}{{\~n}}1 {Ñ}{{\~N}}1 {¿}{{?`}}1 {¡}{{!`}}1
}

% Metadatos institucionales por defecto
\title[Modelación Estadística]{Modelación Estadística para Ciencia de Datos e Ingeniería}
\author[J. Castillo Colmenares]{Juliho David Castillo Colmenares}
\institute[Tec de Monterrey]{Tecnológico de Monterrey \\ Escuela de Ingeniería y Ciencias}
\date{\today}

% Comandos y macros estadísticas compartidas para las presentaciones Beamer (Metropolis)

% Conjuntos numéricos básicos
\newcommand{\R}{\mathbb{R}}
\newcommand{\N}{\mathbb{N}}
\newcommand{\Q}{\mathbb{Q}}
\newcommand{\Z}{\mathbb{Z}}
\newcommand{\set}[1]{\left\{ #1 \right\}}
\newcommand{\abs}[1]{\left|#1\right|}
\newcommand{\norm}[1]{\left\| #1 \right\|}

% Comandos estadísticos y probabilísticos del curso
\newcommand{\Var}{\operatorname{Var}}
\newcommand{\cov}{\operatorname{Cov}}
\newcommand{\corr}{\rho}
\newcommand{\comb}[2]{\begin{pmatrix} #1 \\ #2 \end{pmatrix}}
\newcommand{\s}{\sigma}
\newcommand{\particion}{\mathfrak{P}}
\newcommand{\card}[1]{\#\left( #1 \right)}
\newcommand{\seq}[3]{\set{#1}_{#2}^{#3}}
\newcommand{\E}{\mathbb{E}}
\newcommand{\Prob}{\mathbb{P}}

% Entornos pedagógicos y cajas en español académico natural
\newenvironment{discusion}{%
  \begin{alertblock}{Pregunta de Discusión}%
}{%
  \end{alertblock}%
}

\newenvironment{reflexion}{%
  \begin{alertblock}{Reflexión para el Aula}%
}{%
  \end{alertblock}%
}

\newenvironment{paradoja}{%
  \begin{alertblock}{Paradoja de Exploración}%
}{%
  \end{alertblock}%
}

\newenvironment{motivacion}{%
  \begin{exampleblock}{Motivación: Ciencia de Datos en la Práctica}%
}{%
  \end{exampleblock}%
}

\newenvironment{retoclase}{%
  \begin{block}{Reto Rápido en Equipo}%
}{%
  \end{block}%
}


\title[Set Notation and Partitions]{Section 2.2: Set Notation and Partitions}
\subtitle{The Algebraic and Causal Language of Randomness}
\author{Julio César Hernández Castro}
\institute{Tecnológico de Monterrey}
\date{}

\begin{document}

% Slide 1: Title Page
\begin{frame}[plain]
  \titlepage
\end{frame}

% Slide 2: Roadmap and Position Indicator
\begin{frame}{Roadmap — Chapter 02: Probability Theory}
  Where are we on our journey through probability theory? \pause
  \begin{itemize}
    \item \textbf{02.01 Introduction to Probability} \emph{(Completed)} \pause
    \item \color{TecRojo}\textbf{02.02 Set Notation and Partitions \emph{(Today)}}\color{black} \pause
    \item \textbf{02.03 Fundamentals and Axioms of Probability} \emph{(Next session)} \pause
    \item \textbf{02.04 Conditional Probability and Independence} \pause
    \item \textbf{02.05 Bayes' Theorem} \pause
    \item \textbf{02.06 Random Sampling}
  \end{itemize}
  \vspace{0.6em}
  \begin{block}{Section 2.2 Objective}
    Standardize set operations and use \emph{partitions} as algebraic puzzles to quantify elements and prevent double counting in data-driven causal analysis.
  \end{block}
\end{frame}

% Slide 3: Motivation / The Survey Dilemma
\begin{frame}{The Counting Dilemma in Data Science}
  \begin{motivacion}
    In a survey of $100$ users of a banking app, $60$ use \textbf{iOS} devices, $45$ use \textbf{Android}, and $30$ use both platforms for their financial transactions. If a manager sums $60 + 45 = 105$, it exceeds the total sample size. Why does this counting paradox occur?
  \end{motivacion}
  
  \vspace{0.4em}
  \pause
  \begin{itemize}
    \item The \textbf{double counting error} arises when ignoring the algebraic structure of intersections in sets that are not mutually exclusive. \pause
    \item To measure probabilities with rigorous accuracy, we must decompose the sample space into a \textbf{puzzle of disjoint partitions}.
  \end{itemize}
\end{frame}

% Slide 4: Set Algebra and Basic Operations
\begin{frame}{Fundamental Set Operations}
  Given a universe or sample space $S$ and subsets $A, B \subset S$:
  
  \vspace{0.2em}
  \begin{itemize}
    \item \textbf{Union ($A \cup B$):} Elements in $A$, in $B$, or in both: $\{x \in S \mid x \in A \text{ or } x \in B\}$. \pause
    \item \textbf{Intersection ($A \cap B$):} Elements simultaneously in both: $\{x \in S \mid x \in A \text{ and } x \in B\}$. \pause
    \item \textbf{Difference ($A \setminus B$):} Elements in $A$ not belonging to $B$: $\{x \in S \mid x \in A \text{ and } x \notin B\}$. \pause
    \item \textbf{Complement ($A'$):} Elements of $S$ outside of $A$: $\{x \in S \mid x \notin A\} = S \setminus A$. \pause
    \item \textbf{Symmetric Difference ($A \triangle B$):} Elements in $A$ or $B$, but not in both: $(A \setminus B) \cup (B \setminus A)$.
  \end{itemize}
\end{frame}

% Slide 5: The Partition Concept (Puzzle)
\begin{frame}{What is a Partition of a Set?}
  A \textbf{partition} of a set $S$ is a finite collection of subsets $\{P_i\}_{i=0}^{N}$ that divides the entire space into exact, non-overlapping pieces:
  
  \vspace{0.4em}
  \begin{block}{Formal Definition of a Partition}
    The collection $\{P_i\}_{i=0}^{N}$ is a partition of $S$ if and only if it satisfies two conditions:
    \begin{enumerate}
      \item \textbf{Mutually Disjoint:} $P_i \cap P_j = \emptyset$ whenever $i \neq j$. No puzzle piece overlaps with another. \pause
      \item \textbf{Total Coverage:} $\bigcup_{i=0}^{N} P_i = S$. Joining all pieces reconstructs the complete space without leaving gaps.
    \end{enumerate}
  \end{block}
\end{frame}

% Slide 6: Standard Partition of 2 Sets (4 Regions)
\begin{frame}{Standard Partition for Two Sets}
  For two events $A, B \subset S$, every element either belongs or does not belong to each set ($2^2 = 4$ elementary disjoint regions):
  
  \vspace{0.3em}
  \begin{itemize}
    \item $E_0 = A' \cap B'$ \quad (Elements in neither of the two categories). \pause
    \item $E_1 = A' \cap B$ \quad (Elements exclusive to set $B$). \pause
    \item $E_2 = A \cap B'$ \quad (Elements exclusive to set $A$). \pause
    \item $E_3 = A \cap B$ \quad (Elements in the simultaneous intersection).
  \end{itemize}
  
  \vspace{0.3em}
  \pause
  \begin{block}{Exact Additive Rule on Partitions}
    Since $E_i \cap E_j = \emptyset$, cardinality counting is strictly linear:
    \[ \card{S} = x_0 + x_1 + x_2 + x_3, \quad \text{where } x_k = \card{E_k} \]
  \end{block}
\end{frame}

% Slide 7: Connection with Linear Systems
\begin{frame}{Algebraic Solution via Linear Systems}
  Any counting problem over $A, B$ translates into a linear system of equations over the partition unknowns ($x_0, x_1, x_2, x_3$):
  
  \vspace{0.3em}
  \begin{center}
    \begin{tabular}{l l l}
      \textbf{Set or Condition} & \textbf{Sum over Partition} & \textbf{iOS/Android Example} \\
      \hline
      Sample Space ($S$) & $x_0 + x_1 + x_2 + x_3$ & $= 100$ \\
      Set $A$ (\emph{iOS}) & $x_2 + x_3$ & $= 60$ \\
      Set $B$ (\emph{Android}) & $x_1 + x_3$ & $= 45$ \\
      Intersection ($A \cap B$) & $x_3$ & $= 30$ \\
      \hline
    \end{tabular}
  \end{center}
  
  \vspace{0.3em}
  \pause
  Substituting $x_3 = 30$, we immediately obtain $x_2 = 30$, $x_1 = 15$, and $x_0 = 25$. The exact sum is $30 + 30 + 15 + 25 = 100$!
\end{frame}

% Slide 8: Partition of 3 Sets (8 Regions)
\begin{frame}{Standard Partition for Three Sets}
  Introducing a third set $C \subset S$ divides the space into $2^3 = 8$ elementary pieces, indexable using binary notation:
  
  \vspace{0.2em}
  \begin{columns}[T]
    \begin{column}{0.48\textwidth}
      \small
      \begin{itemize}
        \item $E_0 = A' \cap B' \cap C'$ \quad ($000$)
        \item $E_1 = A' \cap B' \cap C$ \quad ($001$)
        \item $E_2 = A' \cap B \cap C'$ \quad ($010$)
        \item $E_3 = A' \cap B \cap C$ \quad ($011$)
      \end{itemize}
    \end{column}
    \begin{column}{0.48\textwidth}
      \small
      \begin{itemize}
        \item $E_4 = A \cap B' \cap C'$ \quad ($100$)
        \item $E_5 = A \cap B' \cap C$ \quad ($101$)
        \item $E_6 = A \cap B \cap C'$ \quad ($110$)
        \item $E_7 = A \cap B \cap C$ \quad ($111$)
      \end{itemize}
    \end{column}
  \end{columns}
  
  \vspace{0.3em}
  \pause
  \begin{center}
    \small This establishes the fundamental \textbf{Inclusion-Exclusion Formula}:
    \[ \card{A \cup B \cup C} = \sum \card{A_i} - \sum \card{A_i \cap A_j} + \card{A \cap B \cap C} \]
  \end{center}
\end{frame}

% Slide 9A: Python Lab (Part 1: Basic Operations)
\begin{frame}[fragile]{Python Lab: Basic Set Operations}
  The script \texttt{02.02\_conjuntos\_particiones.py} models event algebra using Python's native \texttt{set}:
  \vspace{0.2em}
  {\lstset{basicstyle=\fontsize{6.5pt}{7.5pt}\selectfont\ttfamily}
  \lstinputlisting[firstline=1, lastline=16]{../../code/02_teoria_probabilidad/02.02_conjuntos_particiones.py}}
\end{frame}

% Slide 9B-1: Python Lab (Part 2: System Formulation)
\begin{frame}[fragile]{Python Lab: Partition Equations}
  Setting up the 4 elementary regions as unknowns in a linear algebra system:
  \vspace{0.1em}
  {\lstset{basicstyle=\fontsize{6.5pt}{7.5pt}\selectfont\ttfamily}
  \lstinputlisting[firstline=18, lastline=35]{../../code/02_teoria_probabilidad/02.02_conjuntos_particiones.py}}
\end{frame}

% Slide 9B-2: Python Lab (Part 3: Numpy Matrix)
\begin{frame}[fragile]{Python Lab: Linear System Matrix}
  Matrix definition $M x = b$ using \texttt{numpy.array}:
  \vspace{0.1em}
  {\lstset{basicstyle=\fontsize{6.5pt}{7.5pt}\selectfont\ttfamily}
  \lstinputlisting[firstline=37, lastline=46]{../../code/02_teoria_probabilidad/02.02_conjuntos_particiones.py}}
\end{frame}

% Slide 9C: Python Lab (Part 4: Solution and Verification)
\begin{frame}[fragile]{Python Lab: Solution and Inspection}
  Matrix solution with \texttt{numpy.linalg.solve} and cardinality validation:
  \vspace{0.1em}
  {\lstset{basicstyle=\fontsize{6.5pt}{7.5pt}\selectfont\ttfamily}
  \lstinputlisting[firstline=47, lastline=54]{../../code/02_teoria_probabilidad/02.02_conjuntos_particiones.py}}
\end{frame}

% Slide 9D: Python Lab (Terminal Output)
\begin{frame}[fragile]{Python Lab: Terminal Output}
  Execution console verifying exact partition cardinalities:
  \vspace{0.1em}
  \begin{block}{Standard Terminal Output}
    \fontsize{6.5pt}{7.5pt}\selectfont
    \begin{verbatim}
--- 1. Operaciones Básicas de Conjuntos ---
Espacio Muestral (S):  [1, 2, 3, 4, 5, 6, 7, 8, 9, 10]
Evento A (Pares):      [2, 4, 6, 8, 10]
Evento B (<= 5):       [1, 2, 3, 4, 5]
Unión (A union B):     [1, 2, 3, 4, 5, 6, 8, 10]
Intersección (A & B):  [2, 4]
Diferencia (A \ B):    [6, 8, 10]
Complemento (A'):      [1, 3, 5, 7, 9]

--- 2. Resolución Algebraica vía Partición Estándar (4 Piezas) ---
x3 (|I & F|   - Ambos idiomas):           30 personas
x2 (|I \ F|   - Solo Inglés):            30 personas
x1 (|F \ I|   - Solo Francés):           15 personas
x0 (|I' & F'| - Ninguno de los dos):      25 personas
Total verificado de la partición:        100 personas
    \end{verbatim}
  \end{block}
\end{frame}

% Slide 10: Textbook Exercises (Tiered 3-3-2-2) - Fundamental
\begin{frame}{Textbook Exercises — Level 1: Fundamental}
  \begin{block}{Problem 2.1.2 (International Conference)}
    Out of $100$ attendees at a conference, $60$ speak English ($I$), $45$ speak French ($F$), and $30$ speak both languages ($I \cap F$). Find:
    \begin{enumerate}
      \item $\card{I \cup F}$, the number who speak at least one of the two languages. \pause
      \item $\card{I \setminus F}$, attendees who speak strictly English only. \pause
      \item $\card{(I \cup F)'}$, the number of attendees who speak neither English nor French.
    \end{enumerate}
  \end{block}
  
  \vspace{0.3em}
  \pause
  \begin{alertblock}{Immediate Solution via Partitions}
    With $x_3=30$, $x_2=30$, $x_1=15$, and $x_0=25$:
    1) $\card{I \cup F} = 75$; \quad 2) $\card{I \setminus F} = 30$; \quad 3) $\card{(I \cup F)'} = 25$.
  \end{alertblock}
\end{frame}

% Slide 11: Textbook Exercises (Tiered 3-3-2-2) - Operational
\begin{frame}{Textbook Exercises — Level 2: Operational}
  \begin{block}{Problem 2.1.4 (Extension to Three Subsets)}
    Extend the partition construction to three subsets $A, B, C \subset S$. How many elements does the partition $\particion(A,B,C)$ contain, and what algebraic property guarantees no overlaps?
  \end{block}
  
  \vspace{0.4em}
  \pause
  \begin{itemize}
    \item \textbf{Number of Elements:} Each set doubles the previous regions ($2^3 = 8$ elementary disjoint regions $E_0, E_1, \dots, E_7$). \pause
    \item \textbf{Disjointness Guarantee:} For any $i \neq j$, their binary representations differ in at least one bit (e.g., belonging to $A$ vs. belonging to $A'$), whose intersection is by definition empty ($\emptyset$).
  \end{itemize}
\end{frame}

% Slide 12: Textbook Exercises (Tiered 3-3-2-2) - Analytical
\begin{frame}{Textbook Exercises — Level 3: Analytical}
  \small
  \begin{block}{Problem 2.1.8 (Inclusion-Exclusion for 3 Sets)}
    Prove the cardinality identity for the union of three sets:
    \begin{align*}
      \card{A \cup B \cup C} = \card{A} + \card{B} + \card{C} - \card{A \cap B} - \card{A \cap C} \\
      - \card{B \cap C} + \card{A \cap B \cap C}
    \end{align*}
  \end{block}
  
  \pause
  \footnotesize \textbf{Proof via Partitions ($x_k = \card{E_k}$):}
  \begin{itemize}
    \item $\card{A} + \card{B} + \card{C}$ counts pairwise intersections $2$ times and the triple intersection $3$ times. \pause
    \item Subtracting the pairwise terms balances double counting but eliminates $x_7$ completely ($3 - 3 = 0$). \pause
    \item Therefore, we add $+ \card{A \cap B \cap C} = +x_7$ back at the end for exact, overlap-free coverage.
  \end{itemize}
\end{frame}

% Slide 13: Textbook Exercises (Tiered 3-3-2-2) - Challenging
\begin{frame}{Textbook Exercises — Level 4: Challenging}
  \small
  \begin{block}{Problem 2.1.10 (Census Absurdities in Epidemiology)}
    For $500$ patients, a report states: $\card{S_1}=300, \card{S_2}=250, \card{S_3}=200$; $\card{S_1 \cap S_2}=180, \card{S_1 \cap S_3}=150, \card{S_2 \cap S_3}=130$; $\card{S_1 \cap S_2 \cap S_3}=30$. Why is this mathematically false?
  \end{block}
  
  \pause
  \begin{alertblock}{Analysis of the Complement Region ($x_0$)}
    Simple pairwise intersections: $\card{S_1 \cap S_2 \cap S_3'} = 150$, $\card{S_1 \cap S_3 \cap S_2'} = 120$, $\card{S_2 \cap S_3 \cap S_1'} = 100$.
    \pause
    Elementary sum: $150+120+100+30 + \text{simple regions} > 500$. Thus $x_0 < 0$, which contradicts non-negativity.
  \end{alertblock}
\end{frame}

% Slide 14: Synthesis and Next Step
\begin{frame}{Summary of Section 2.2}
  \begin{block}{Key Takeaways}
    \begin{itemize}
      \item Any operation involving non-exclusive sets must be transformed into a \textbf{disjoint partition} to guarantee additive, bias-free sums. \pause
      \item Linear systems of equations over partitions ($Mx = b$) are the optimal computational tool to validate surveys and census data.
    \end{itemize}
  \end{block}
  
  \vspace{0.6em}
  \pause
  {\large \color{TecAzulOscuro}\textbf{Next Session: 02.03 Fundamentals and Axioms of Probability}}
  \vspace{0.3em}
  \begin{itemize}
    \item With set notation mastered, we will introduce the probability measure function $P(A)$ and Kolmogorov's three immortal axioms.
  \end{itemize}
\end{frame}

\end{document}
